\documentclass[12pt,a4paper]{report}
\usepackage[utf8]{inputenc}
\usepackage{amsmath}
\usepackage{amsfonts}
\usepackage{amssymb}
\usepackage{amsthm}
\usepackage{hyperref}

\usepackage{multicol}
\usepackage{fancyhdr}
\usepackage[inline]{enumitem}
\usepackage{tikz}
\usepackage{tikz-cd}
\usetikzlibrary{calc}
\usetikzlibrary{shapes.geometric}
\usepackage[margin=0.5in]{geometry}
\usepackage{xcolor}
\usepackage{ esint }

\hypersetup{
    colorlinks=true,
    linkcolor=blue,
    filecolor=magenta,      
    urlcolor=cyan,
    pdftitle={PDF2 Top 10},
    pdfpagemode=FullScreen,
    }

%\urlstyle{same}

\newcommand{\CLASSNAME}{Comprehensive Exam -- PDE2, Top 10}
\newcommand{\STUDENTNAME}{Paul Carmody}
\newcommand{\ASSIGNMENT}{NOTES }
\newcommand{\DUEDATE}{April 29, 2026}
\newcommand{\SEMESTER}{Spring 2026}
\newcommand{\SCHEDULE}{3:00 -- 5:00}
\newcommand{\ROOM}{Crawford 332}

\newcommand{\MMN}{M_{m\times n}}
\newcommand{\FF}{\mathcal{F}}

\pagestyle{fancy}
\fancyhf{}
\chead{ \fancyplain{}{\CLASSNAME} }
%\chead{ \fancyplain{}{\STUDENTNAME} }
\rhead{\thepage}
\fancyfoot[C]{-- \thepage --}
\input{../../MyMacros}

\newcommand{\RED}[1]{\textcolor{red}{#1}}
\newcommand{\BLUE}[1]{\textcolor{blue}{#1}}
\newcommand{\FFF}{\mathcal{F}}

\begin{document}

\begin{center}
	\Large{\CLASSNAME -- \SEMESTER} \\
	\large{ w/Faculty}
\end{center}
\begin{center}
	\STUDENTNAME \\
	\ASSIGNMENT -- \DUEDATE\\
\end{center} 


Here is a \DEFINE{Top 10 PDE Comprehensive Exam Practice Set}, designed around the kinds of problems that repeatedly appear on graduate qualifying/comprehensive exams.

These focus on the ``core moves" of graduate PDE:

	\begin{itemize}
    \item weak formulations,
    \item energy estimates,
    \item Sobolev embeddings,
    \item compactness,
    \item maximum principles,
    \item Fourier methods,
    \item variational methods.
	\end{itemize}

These are the backbone of graduate Partial Differential Equations.
\begin{description}

\newpage\item \textbf{1. Energy estimate for the heat equation}

Let $u$ solve
\begin{align*}
u_t-\Delta u=0 \quad \text{in }U\times(0,T),
\end{align*}
with $u=0$ on $\partial U\times(0,T)$.

Show that
\begin{align*}
\frac{d}{dt}\int_U u^2\,dx\le 0.
\end{align*}
Then prove uniqueness.

Skills tested:

	\begin{itemize}
    \item multiply by $u$,
    \item integration by parts,
    \item energy decay.
	\end{itemize}

\HLINE


We use the standard \DEFINE{energy method} for the heat equation.

Let $u$ solve
\begin{align*}
u_t-\Delta u=0
\qquad\text{in }U\times(0,T),
\end{align*}
with boundary condition
\begin{align*}
u=0
\qquad\text{on }\partial U\times(0,T).
\end{align*}
We will show:

\begin{align*}
\frac{d}{dt}\int_U u^2\,dx\le 0,
\end{align*}
and then use this to prove uniqueness.

This is the basic energy estimate for the heat equation in Partial Differential Equations.

\begin{description}
\item \textbf{Step 1: Multiply the PDE by $u$}

Starting with
\begin{align*}
u_t-\Delta u=0,
\end{align*}
multiply by $u$:
\begin{align*}
uu_t-u\Delta u=0.
\end{align*}
Integrate over $U$:
\begin{align*}
\int_U uu_t\,dx-\int_U u\Delta u\,dx=0.
\end{align*}

\item \textbf{Step 2: Integrate by parts}

Using integration by parts,
\begin{align*}
\int_U u\Delta u\,dx = -\int_U |\nabla u|^2\,dx + \int_{\partial U}u\frac{\partial u}{\partial\nu}\,dS.
\end{align*}
Since $u=0$ on $\partial U$, the boundary term vanishes:
\begin{align*}
\int_{\partial U}u\frac{\partial u}{\partial\nu}\,dS=0.
\end{align*}
Thus,
\begin{align*}
\int_U u\Delta u\,dx = -\int_U |\nabla u|^2\,dx.
\end{align*}
Substitute back:
\begin{align*}
\int_U uu_t\,dx+\int_U |\nabla u|^2\,dx=0.
\end{align*}

\item \textbf{Step 3: Recognize the energy derivative}

Since
\begin{align*}
\frac{d}{dt}\int_U u^2\,dx = 2\int_U uu_t\,dx,
\end{align*}
we obtain
\begin{align*}
\frac12\frac{d}{dt}\int_U u^2\,dx+\int_U |\nabla u|^2\,dx=0.
\end{align*}
Therefore,
\begin{align*}
\frac{d}{dt}\int_U u^2\,dx = -2\int_U |\nabla u|^2\,dx.
\end{align*}
Hence,
\begin{align*}
\boxed{
\frac{d}{dt}\int_U u^2\,dx\le 0.
}
\end{align*}
So the $L^2$-energy is nonincreasing.

\item \textbf{Step 4: Prove uniqueness}

Suppose $u_1,u_2$ are two solutions with the same initial data.

Define
\begin{align*}
w=u_1-u_2.
\end{align*}
Then $w$ satisfies
\begin{align*}
w_t-\Delta w=0
\end{align*}
in $U\times(0,T)$, with
\begin{align*}
w=0
\end{align*}
on $\partial U\times(0,T)$, and initial condition
\begin{align*}
w(x,0)=0.
\end{align*}
Apply the energy identity to $w$:
\begin{align*}
\frac{d}{dt}\int_U w^2\,dx\le 0.
\end{align*}
But initially,
\begin{align*}
\int_U w(x,0)^2\,dx=0.
\end{align*}
Since the energy is nonincreasing and nonnegative,
\begin{align*}
\int_U w(x,t)^2\,dx=0
\qquad\forall t.
\end{align*}
Thus,
\begin{align*}
w(x,t)=0
\qquad\forall (x,t).
\end{align*}
Therefore,
\begin{align*}
u_1=u_2.
\end{align*}
Hence the solution is unique.

\item \textbf{Final result}

We proved:
\begin{align*}
\boxed{
\frac{d}{dt}\int_U u^2\,dx = -2\int_U |\nabla u|^2\,dx \le 0 }
\end{align*}
and this implies:
\begin{align*}
\boxed{\text{The solution of the heat equation is unique.}}
\end{align*}

\item \textbf{Interpretation}

The quantity
\begin{align*}
\int_U u^2\,dx
\end{align*}
is the heat energy, and diffusion causes it to decrease over time.

That decay prevents two different solutions from starting from the same initial data.

\end{description}

If you'd like, next I can show how this energy estimate yields the \DEFINE{stability estimate}

\begin{align*}
|u(\cdot,t)|*{L^2(U)}\le |g|*{L^2(U)},
\end{align*}

which is another common comprehensive exam result.


\newpage\item \textbf{2. Energy conservation for the wave equation}

Let $u$ solve

\begin{align*}
u_{tt}-\Delta u=0
\end{align*}
with homogeneous Dirichlet boundary data.

Show that
\begin{align*}
E(t)=\frac12\int_U(u_t^2+|\nabla u|^2)\,dx
\end{align*}
is constant.

Skills tested:

	\begin{itemize}
    \item multiply by $u_t$,
    \item conserved energy,
    \item uniqueness.
	\end{itemize}

\HLINE


We prove that the energy
\begin{align*}
E(t)=\frac12\int_U\left(u_t^2+|\nabla u|^2\right)\,dx
\end{align*}
is constant for solutions of
\begin{align*}
u_{tt}-\Delta u=0
\qquad\text{in }U\times(0,T),
\end{align*}
with homogeneous Dirichlet boundary condition
\begin{align*}
u=0 \qquad\text{on }\partial U\times(0,T).
\end{align*}

This is the standard \DEFINE{energy conservation law} for the wave equation in Partial Differential Equations.
\begin{description}

\item \textbf{Step 1: Differentiate the energy}

Differentiate $E(t)$:
\begin{align*}
E'(t) = \frac12\frac{d}{dt}\int_U\left(u_t^2+|\nabla u|^2\right)\,dx.
\end{align*}
So
\begin{align*}
E'(t) = \int_U u_tu_{tt}\,dx + \int_U \nabla u\cdot\nabla u_t\,dx.
\end{align*}

\item \textbf{Step 2: Integrate by parts}

Integrate the second term by parts:
\begin{align*}
\int_U \nabla u\cdot\nabla u_t\,dx = -\int_U (\Delta u)u_t\,dx + \int_{\partial U}\frac{\partial u}{\partial\nu}u_t\,dS.
\end{align*}
Thus,
\begin{align*}
E'(t) =  \int_U u_tu_{tt}\,dx- \int_U (\Delta u)u_t\,dx + \int_{\partial U}\frac{\partial u}{\partial\nu}u_t\,dS.
\end{align*}
Factor $u_t$:
\begin{align*}
E'(t) = \int_U u_t(u_{tt}-\Delta u)\,dx + \int_{\partial U}\frac{\partial u}{\partial\nu}u_t\,dS.
\end{align*}

\item \textbf{Step 3: Use the PDE}

Since
\begin{align*}
u_{tt}-\Delta u=0,
\end{align*}
the volume term vanishes:
\begin{align*}
\int_U u_t(u_{tt}-\Delta u)\,dx=0.
\end{align*}
Hence
\begin{align*}
E'(t) = \int_{\partial U}\frac{\partial u}{\partial\nu}u_t\,dS.
\end{align*}

\item \textbf{Step 4: Use the boundary condition}

Since
\begin{align*}
u=0 \qquad\text{on }\partial U\times(0,T),
\end{align*}
we have
\begin{align*}
u_t=0
\qquad\text{on }\partial U\times(0,T),
\end{align*}
because the boundary value is identically zero in time.

Thus,
\begin{align*}
\int_{\partial U}\frac{\partial u}{\partial\nu}u_t\,dS=0.
\end{align*}
Therefore,
\begin{align*}
E'(t)=0.
\end{align*}

\item \textbf{Step 5: Conclude conservation}

Since
\begin{align*}
E'(t)=0,
\end{align*}

the energy is constant:
\begin{align*}
\boxed{
E(t)=\frac12\int_U\left(u_t^2+|\nabla u|^2\right)\,dx
\text{ is constant in time.}
}
\end{align*}

\item \textbf{Final conclusion}

For the wave equation with homogeneous Dirichlet boundary data,
\begin{align*}
\boxed{
\frac{d}{dt} \left[ \frac12\int_U (u_t^2+|\nabla u|^2)\,dx \right] =0
}
\end{align*}
so the total energy is conserved.

\item \textbf{Interpretation}

The two terms represent:
\begin{align*}
\frac12\int_U u_t^2\,dx
\quad\text{(kinetic energy)}
\end{align*}
and
\begin{align*}
\frac12\int_U |\nabla u|^2\,dx
\quad\text{(potential energy)}.
\end{align*}
The wave equation transfers energy between these two forms, but the total remains constant.

\end{description}
---

If you'd like, next I can show how this same computation gives \DEFINE{uniqueness of the wave equation}, exactly as the heat equation energy estimate gave uniqueness.


\newpage\item \textbf{3. Weak formulation of Poisson’s equation}

Given
\begin{align*}
-\Delta u=f \quad \text{in }U,\qquad u=0 \text{ on }\partial U,
\end{align*}
derive the weak formulation:
\begin{align*}
\int_U \nabla u\cdot\nabla v\,dx=\int_U fv\,dx.
\end{align*}

Skills tested:
	\begin{itemize}
    \item integration by parts,
    \item Sobolev spaces,
    \item weak derivatives.
	\end{itemize}

\HLINE


To derive the \DEFINE{weak formulation} of
\begin{align*}
-\Delta u=f
\quad\text{in }U,
\qquad
u=0
\quad\text{on }\partial U,
\end{align*}
we convert the PDE into an integral identity using \DEFINE{test functions} and \DEFINE{integration by parts}.

This is the standard weak formulation for Poisson’s equation in Partial Differential Equations.

\begin{description}

\item \textbf{Step 1: Choose a test function}

Let
\begin{align*}
v\in C_c^\infty(U),
\end{align*}
that is, $v$ is smooth and compactly supported in $U$.

Multiply the PDE by $v$:
\begin{align*}
-\Delta u\, v=fv.
\end{align*}
Integrate over $U$:
\begin{align*}
\int_U (-\Delta u)v\,dx=\int_U fv\,dx.
\end{align*}

\item \textbf{Step 2: Integrate by parts}

Use Green’s identity:
\begin{align*}
\int_U (-\Delta u)v\,dx =  \int_U \nabla u\cdot\nabla v\,dx - \int_{\partial U}v\frac{\partial u}{\partial\nu}\,dS.
\end{align*}
Since $v\in C_c^\infty(U$), we have
\begin{align*}
v=0 \quad\text{on }\partial U,
\end{align*}
so the boundary term vanishes:
\begin{align*}
\int_{\partial U}v\frac{\partial u}{\partial\nu}\,dS=0.
\end{align*}
Thus,
\begin{align*}
\int_U (-\Delta u)v\,dx = \int_U \nabla u\cdot\nabla v\,dx.
\end{align*}
Substituting into the previous equation gives:
\begin{align*}
\int_U \nabla u\cdot\nabla v\,dx = \int_U fv\,dx.
\end{align*}

\item \textbf{Step 3: Extend to Sobolev test functions}

The above identity holds for all
\begin{align*}
v\in C_c^\infty(U).
\end{align*}
Since $C_c^\infty(U)$ is dense in $H_0^1(U)$, we extend the identity to all
\begin{align*}
v\in H_0^1(U).
\end{align*}
Thus the weak formulation is:
\begin{align*}
\boxed{
\int_U \nabla u\cdot\nabla v\,dx = \int_U fv\,dx \qquad\forall v\in H_0^1(U).
}
\end{align*}

\item \textbf{Final answer}

The weak formulation of
\begin{align*}
-\Delta u=f,\qquad u=0\text{ on }\partial U
\end{align*}
is:
\begin{align*}
\boxed{
\int_U \nabla u\cdot\nabla v\,dx = \int_U fv\,dx \qquad\forall v\in H_0^1(U).
}
\end{align*}

\item \textbf{Why this is useful}

The classical PDE requires second derivatives:
\begin{align*}
\Delta u.
\end{align*}
But the weak formulation only requires first derivatives:
\begin{align*}
\nabla u.
\end{align*}
So instead of needing $u\in C^2(U)$, we only need
\begin{align*}
u\in H_0^1(U),
\end{align*}
which makes existence theory much easier.

\end{description}

If you'd like, next I can show how to use this weak formulation with \DEFINE{Lax–Milgram} to prove existence and uniqueness of $u$.


\newpage\item \textbf{4. Existence by Lax–Milgram}

Show there exists a unique weak solution to

\begin{align*}
-\Delta u+u=f,\qquad u|_{\partial U}=0
\end{align*}

for $f\in L^2(U)$.

Skills tested:

	\begin{itemize}
    \item bilinear form,
    \item boundedness,
    \item coercivity,
    \item Lax–Milgram.
	\end{itemize}

\HLINE


We will prove existence and uniqueness of a weak solution using the \DEFINE{Lax–Milgram theorem}, the standard method for elliptic problems in Partial Differential Equations.

\begin{description}

\item \textbf{Step 1: Write the weak formulation}

We consider
\begin{align*}
-\Delta u+u=f \quad\text{in }U, \qquad u=0 \quad\text{on }\partial U.
\end{align*}
Multiply by a test function $v\in H_0^1(U)$:
\begin{align*}
(-\Delta u+u)v=fv.
\end{align*}
Integrate over $U$:
\begin{align*}
\int_U(-\Delta u)v\,dx+\int_U uv\,dx=\int_U fv\,dx.
\end{align*}
Integrate the Laplacian term by parts:
\begin{align*}
\int_U(-\Delta u)v\,dx=\int_U \nabla u\cdot\nabla v\,dx,
\end{align*}
since $v=0$ on $\partial U$.

Thus the weak formulation is:
\begin{align*}
\boxed{
\int_U (\nabla u\cdot\nabla v+uv)\,dx = \int_U fv\,dx \qquad\forall v\in H_0^1(U).
}
\end{align*}

\item \textbf{Step 2: Define the bilinear form}

Define
\begin{align*}
a(u,v)=\int_U(\nabla u\cdot\nabla v+uv)\,dx.
\end{align*}
Define the linear functional
\begin{align*}
L(v)=\int_U fv\,dx.
\end{align*}
Then the weak problem is:
\begin{align*}
a(u,v)=L(v)
\qquad\forall v\in H_0^1(U).
\end{align*}
Now we verify the hypotheses of \DEFINE{Lax–Milgram}.

\item \textbf{Step 3: Show $a(\cdot,\cdot)$ is bounded}

We estimate:
\begin{align*}
|a(u,v)| \le \int_U |\nabla u||\nabla v|\,dx+\int_U |u||v|\,dx.
\end{align*}
Using Cauchy–Schwarz:
\begin{align*}
|a(u,v)| \le \|\nabla u\|_{L^2}\|\nabla v\|_{L^2} + \|u\|_{L^2}\|v\|_{L^2}.
\end{align*}
Hence,
\begin{align*}
|a(u,v)| \le \|u\|_{H_0^1(U)}\|v\|_{H_0^1(U)}.
\end{align*}
So $a$ is bounded.

\item \textbf{Step 4: Show coercivity}

Compute:
\begin{align*}
a(u,u)=\int_U (|\nabla u|^2+u^2)\,dx.
\end{align*}
Thus,
\begin{align*}
a(u,u)=\|u\|_{H_0^1(U)}^2.
\end{align*}
Hence
\begin{align*}
a(u,u)\ge \|u\|_{H_0^1(U)}^2.
\end{align*}
So $a$ is coercive.

\item \textbf{Step 5: Show $L$ is bounded}

We estimate:
\begin{align*}
|L(v)| = \left|\int_U fv\,dx\right| \le \|f\|_{L^2}\|v\|_{L^2}.
\end{align*}
Since
\begin{align*}
\|v\|_{L^2}\le \|v\|_{H_0^1(U)},
\end{align*}
we get
\begin{align*}
|L(v)| \le
\|f\|_{L^2}\|v\|_{H_0^1(U)}.
\end{align*}
So $L$ is a bounded linear functional.

\item \textbf{Step 6: Apply Lax–Milgram}

Since:

\begin{enumerate}
\item $a(u,v$) is bounded,
\item $a(u,u$) is coercive,
\item $L(v$) is bounded,
\end{enumerate}

the \DEFINE{Lax–Milgram theorem} gives a unique $u\in H_0^1(U)$ such that
\begin{align*}
a(u,v)=L(v) \qquad\forall v\in H_0^1(U).
\end{align*}
That is,
\begin{align*}
\int_U (\nabla u\cdot\nabla v+uv)\,dx = \int_U fv\,dx \qquad\forall v\in H_0^1(U).
\end{align*}
So there exists a unique weak solution.

\item \textbf{Final conclusion}

There exists a unique
\begin{align*}
u\in H_0^1(U)
\end{align*}
such that
\begin{align*}
\boxed{
\int_U (\nabla u\cdot\nabla v+uv)\,dx = \int_U fv\,dx \quad\forall v\in H_0^1(U),
}
\end{align*}
and this is the unique weak solution of
\begin{align*}
-\Delta u+u=f,\qquad u|_{\partial U}=0.
\end{align*}

\item \textbf{Key structure}

The proof is always:

\begin{enumerate}
\item derive weak form,
\item define bilinear form,
\item prove boundedness,
\item prove coercivity,
\item apply Lax–Milgram.
\end{enumerate}

This is the standard existence/uniqueness framework for linear elliptic PDE.

\end{description}

If you'd like, next I can show how to prove the \DEFINE{a priori estimate}
\begin{align*}
\|u\|_{H_0^1(U)}\le C\|f\|_{L^2(U)},
\end{align*}
which follows immediately from the same argument and is often asked on exams.


\newpage\item \textbf{5. Sobolev embedding estimate}

Show that if $U\subset\mathbb R^n$ is bounded and $n>2$,

\begin{align*}
H_0^1(U)\hookrightarrow L^{2^*}(U),
\qquad 2^*=\frac{2n}{n-2}.
\end{align*}

Skills tested:

	\begin{itemize}
    \item Sobolev inequality,
    \item critical exponent.
	\end{itemize}

\HLINE


We want to prove the \DEFINE{Sobolev embedding}
\begin{align*}
H_0^1(U)\hookrightarrow L^{2^*}(U),
\qquad 2^*=\frac{2n}{n-2},
\end{align*}
for a bounded domain $U\subset\mathbb R^n$, $n>2$.

That means there exists a constant $C>0$ such that
\begin{align*}
\|u\|_{L^{2^*}(U)} \le C\|u\|_{H_0^1(U)} \qquad\forall u\in H_0^1(U).
\end{align*}
This is the classical \DEFINE{Sobolev inequality}, one of the central results in Functional Analysis.

\begin{description}

\item \textbf{Step 1: Sobolev inequality on $\mathbb R^n$}

The key result is:
If $u\in C_c^\infty(\mathbb R^n$), then
\begin{align*}
\|u\|_{L^{2^*}(\mathbb R^n)} \le C\|\nabla u\|_{L^2(\mathbb R^n)},
\end{align*}
where
\begin{align*}
2^*=\frac{2n}{n-2}.
\end{align*}
This is the Sobolev inequality on $\mathbb R^n$.

\item \textbf{Step 2: Extend functions by zero}

Take $u\in H_0^1(U)$.

Extend $u$ by zero outside $U$:
\begin{align*}
\tilde u(x)=
\begin{cases}
u(x), & x\in U,\\
0, & x\notin U.
\end{cases}
\end{align*}
Then
\begin{align*}
\tilde u\in H^1(\mathbb R^n),
\end{align*}
and
\begin{align*}
\|\tilde u\|_{L^{2^*}(\mathbb R^n)} = \|u\|_{L^{2^*}(U)}.
\end{align*}
Also,
\begin{align*}
\|\nabla \tilde u\|_{L^2(\mathbb R^n)} = \|\nabla u\|_{L^2(U)}.
\end{align*}
Applying the Sobolev inequality in $\mathbb R^n$,
\begin{align*}
\|u\|_{L^{2^*}(U)} \le C\|\nabla u\|_{L^2(U)}.
\end{align*}

\item \textbf{Step 3: Relate to the $H_0^1$-norm}

For $H_0^1(U)$, the norm
\begin{align*}
\|u\|_{H_0^1(U)}
\end{align*}
is equivalent to
\begin{align*}
\|\nabla u\|_{L^2(U)}
\end{align*}
by the \DEFINE{Poincaré inequality}:
\begin{align*}
\|u\|_{L^2(U)} \le C\|\nabla u\|_{L^2(U)}.
\end{align*}
Thus,
\begin{align*}
\|u\|_{H_0^1(U)} \sim \|\nabla u\|_{L^2(U)}.
\end{align*}
Hence,
\begin{align*}
\|u\|_{L^{2^*}(U)} \le C\|u\|_{H_0^1(U)}.
\end{align*}

This proves the embedding.

\item \textbf{Final result}

We have shown that for all $u\in H_0^1(U)$,

\begin{align*}
\boxed{
\|u\|_{L^{2^*}(U)} \le C\|u\|_{H_0^1(U)}, \qquad 2^*=\frac{2n}{n-2}.
}
\end{align*}
Therefore,
\begin{align*}
\boxed{
H_0^1(U)\hookrightarrow L^{2^*}(U).
}
\end{align*}

\item \textbf{Why $2^*=\frac{2n}{n-2}$?}

This exponent is the \DEFINE{critical Sobolev exponent}.
It is the largest exponent for which first derivatives in $L^2$ force the function itself to lie in $L^p$.

The relation comes from scaling:
\begin{align*}
\frac1{2^*}=\frac12-\frac1n.
\end{align*}
That gives
\begin{align*}
2^*=\frac{2n}{n-2}.
\end{align*}

\item \textbf{Key exam takeaway}

To prove Sobolev embedding on a bounded domain:

\begin{enumerate}
\item extend $u$ by zero,
\item apply Sobolev inequality on $\mathbb R^n$,
\item use Poincaré inequality.
\end{enumerate}

This is the standard graduate PDE proof structure.

\end{description}

If you'd like, next I can show the \DEFINE{compact embedding}
\begin{align*}
H_0^1(U)\hookrightarrow\hookrightarrow L^2(U),
\end{align*}
which is the Rellich–Kondrachov theorem and appears constantly on qualifying exams.


\newpage\item \textbf{6. Weak compactness + passing to the limit}

Suppose $u_j$ solves
\begin{align*}
-\Delta u_j=f_j,\qquad u_j\in H_0^1(U),
\end{align*}
where $f_j\rightharpoonup f$ in $L^2(U)$, and $u_j$ is bounded in $H_0^1(U)$.

Show a subsequence converges weakly to a weak solution of
\begin{align*}
-\Delta u=f.
\end{align*}

Skills tested:

	\begin{itemize}
    \item weak convergence,
    \item compactness,
    \item limit in weak formulation.
	\end{itemize}

\HLINE


We are given:

\begin{itemize}
    \item $u_j\in H_0^1(U)$ solves
\begin{align*}
-\Delta u_j=f_j,
\end{align*}

    \item $f_j\rightharpoonup f$ weakly in $L^2(U)$,
    \item $(u_j$) is bounded in $H_0^1(U)$.
\end{itemize}

We want to show that some subsequence $u_{j_k}$ converges weakly in $H_0^1(U)$ to a weak solution $u$ of
\begin{align*}
-\Delta u=f.
\end{align*}
This is a standard \DEFINE{weak compactness + passage to the limit} argument in Functional Analysis.

\begin{description}

\item \textbf{Step 1: Use boundedness in $H_0^1(U)$}

Since $(u_j)$ is bounded in $H_0^1(U)$, and $H_0^1(U)$ is a Hilbert space, bounded sets are weakly precompact.

Therefore there exists a subsequence $u_{j_k}$ and some $u\in H_0^1(U)$ such that
\begin{align*}
u_{j_k}\rightharpoonup u
\qquad\text{weakly in }H_0^1(U).
\end{align*}
That means:
\begin{align*}
\int_U \nabla u_{j_k}\cdot\nabla v\,dx
\to
\int_U \nabla u\cdot\nabla v\,dx
\end{align*}
for every $v\in H_0^1(U)$.

\item \textbf{Step 2: Write the weak formulation for $u_j$}

Since $u_j$ solves
\begin{align*}
-\Delta u_j=f_j,
\end{align*}
its weak formulation is:
\begin{align*}
\int_U \nabla u_j\cdot\nabla v\,dx = \int_U f_jv\,dx \qquad\forall v\in H_0^1(U).
\end{align*}
For the subsequence:
\begin{align*}
\int_U \nabla u_{j_k}\cdot\nabla v\,dx = \int_U f_{j_k}v\,dx.
\end{align*}

\item \textbf{Step 3: Pass to the limit on the left-hand side}

Since
\begin{align*}
u_{j_k}\rightharpoonup u
\quad\text{in }H_0^1(U),
\end{align*}
we have
\begin{align*}
\int_U \nabla u_{j_k}\cdot\nabla v\,dx \to \int_U \nabla u\cdot\nabla v\,dx.
\end{align*}

\item \textbf{Step 4: Pass to the limit on the right-hand side}

Since
\begin{align*}
f_{j_k}\rightharpoonup f
\quad\text{in }L^2(U),
\end{align*}
we know
\begin{align*}
\int_U f_{j_k}\phi\,dx \to \int_U f\phi\,dx \qquad\forall \phi\in L^2(U).
\end{align*}

Because $v\in H_0^1(U)\subset L^2(U)$, we may take $\phi=v$, giving:

\begin{align*}
\int_U f_{j_k}v\,dx \to \int_U fv\,dx.
\end{align*}

\item \textbf{Step 5: Identify the limit equation}

Taking limits in
\begin{align*}
\int_U \nabla u_{j_k}\cdot\nabla v\,dx = \int_U f_{j_k}v\,dx,
\end{align*}
we obtain
\begin{align*}
\int_U \nabla u\cdot\nabla v\,dx = \int_U fv\,dx \qquad\forall v\in H_0^1(U).
\end{align*}
This is exactly the weak formulation of

\begin{align*}
-\Delta u=f,\qquad u\in H_0^1(U).
\end{align*}
Hence $u$ is a weak solution.

\item \textbf{Final conclusion}

There exists a subsequence $u_{j_k}$ and a function $u\in H_0^1(U)$ such that
\begin{align*}
u_{j_k}\rightharpoonup u
\quad\text{weakly in }H_0^1(U),
\end{align*}
and $u$ satisfies
\begin{align*}
\boxed{
\int_U \nabla u\cdot\nabla v\,dx = \int_U fv\,dx \qquad\forall v\in H_0^1(U),
}
\end{align*}
that is,
\begin{align*}
\boxed{-\Delta u=f}
\end{align*}
in the weak sense.

\item \textbf{Key exam structure}

This is the standard 3-step compactness argument:

	\begin{description}
        \item 1. Bounded sequence in $H_0^1$
\begin{align*}
u_j \text{ bounded}
\Rightarrow
u_{j_k}\rightharpoonup u
\end{align*}

        \item 2. Write weak formulation
\begin{align*}
a(u_j,v)=L_j(v)
\end{align*}

        \item 3. Pass to the limit
Use weak convergence on both sides to get
\begin{align*}
a(u,v)=L(v)
\end{align*}

		\end{description}

This exact structure appears constantly in graduate PDE exams.
\end{description}

If you'd like, next I can show how to strengthen this result to \DEFINE{strong convergence in $L^2(U)$} using the \DEFINE{Rellich compactness theorem}, which is often the next step on comprehensive exams.


\newpage\item \textbf{7. Maximum principle}

Let $u$ solve
\begin{align*}
u_t-\Delta u=0
\end{align*}
with $u(x,0)=g(x)$.

Show:
\begin{align*}
\sup_{U\times[0,T]}u\le \sup_U g.
\end{align*}

Skills tested:
	\begin{itemize}
    \item parabolic maximum principle,
    \item uniqueness,
    \item $L^\infty$ bounds.
	\end{itemize}

\HLINE


We want to prove the \DEFINE{maximum principle} for the heat equation:

If $u$ solves

\begin{align*}
u_t-\Delta u=0
\qquad\text{in }U\times(0,T),
\end{align*}
with initial data
\begin{align*}
u(x,0)=g(x),
\end{align*}
then
\begin{align*}
\sup_{U\times[0,T]}u\le \sup_U g.
\end{align*}
This is the standard \DEFINE{parabolic maximum principle}, one of the central tools in Partial Differential Equations.

\begin{description}

\item \textbf{Step 1: Define the maximum level}

Let
\begin{align*}
M=\sup_U g.
\end{align*}
We want to show:
\begin{align*}
u(x,t)\le M
\qquad\forall (x,t)\in U\times[0,T].
\end{align*}
Define
\begin{align*}
v(x,t)=u(x,t)-M.
\end{align*}
Then $v$ satisfies
\begin{align*}
v_t-\Delta v=0,
\end{align*}
and
\begin{align*}
v(x,0)=g(x)-M\le 0.
\end{align*}
So initially, $v\le 0$.

We must show $v\le 0$ for all time.

\item \textbf{Step 2: Suppose $v$ becomes positive}

Assume for contradiction that $v>0$ somewhere in $U\times(0,T]$.

Then $v$ attains a positive maximum at some point
\begin{align*}
(x_0,t_0)\in U\times(0,T].
\end{align*}
At this interior maximum:
\begin{align*}
\nabla v(x_0,t_0)&=0\\
\Delta v(x_0,t_0)&\le 0,
\end{align*}
and since $t_0$ is the first time of maximum,
\begin{align*}
v_t(x_0,t_0)\ge 0.
\end{align*}
Thus,
\begin{align*}
v_t-\Delta v\ge 0.
\end{align*}
But since $v$ satisfies the heat equation,
\begin{align*}
v_t-\Delta v=0.
\end{align*}
Hence at the maximum point,
\begin{align*}
0=v_t-\Delta v\ge 0.
\end{align*}
This forces equality:
\begin{align*}
v_t=0,\qquad \Delta v=0.
\end{align*}
By the strong maximum principle, if a positive interior maximum occurs, then $v$ must be constant.

But $v(x,0)\le 0$, so $v$ cannot be a positive constant.

Contradiction.

Thus $v\le 0$.

\item \textbf{Step 3: Conclude the estimate}

Since $v=u-M\le 0$,
\begin{align*}
u(x,t)\le M.
\end{align*}
Therefore,
\begin{align*}
u(x,t)\le \sup_U g
\qquad\forall (x,t)\in U\times[0,T].
\end{align*}
Hence
\begin{align*}
\boxed{
\sup_{U\times[0,T]}u\le \sup_U g.
}
\end{align*}

\item \textbf{Final result}

We have shown the \DEFINE{maximum principle}:

\begin{align*}
\boxed{
\sup_{U\times[0,T]}u \le \sup_U g.
}
\end{align*}

This means the heat equation cannot create temperatures larger than the initial maximum.

\item \textbf{Interpretation}

For the heat equation,
\begin{align*}
u_t-\Delta u=0,
\end{align*}
heat diffuses from hot regions to cold ones.

So:

	\begin{itemize}
    \item peaks flatten out,
    \item no new maxima are created,
    \item the maximum temperature decreases or stays constant.
	\end{itemize}

That is exactly what the maximum principle says.

\item \textbf{Exam insight}

For heat equation maximum principle problems, the standard strategy is:

\begin{enumerate}
\item let $M=\sup g$,
\item define $v=u-M$,
\item show $v\le 0$ by contradiction using an interior maximum.
\end{enumerate}

This is the standard maximum principle template for parabolic PDE.
\end{description}

If you'd like, next I can show how the same argument proves the \DEFINE{minimum principle}
\begin{align*}
\inf_{U\times[0,T]}u\ge \inf_U g,
\end{align*}
and combining them gives the useful estimate
\begin{align*}
\|u(\cdot,t)\|_{L^\infty(U)}\le \|g\|_{L^\infty(U)}.
\end{align*}


\newpage\item \textbf{8. Direct method in the calculus of variations}

Show that the infimum
\begin{align*}
\inf_{u\in H_0^1(U)} \left( \int_U |\nabla u|^2\,dx-\int_U fu\,dx \right)
\end{align*}
is attained.

Skills tested:

	\begin{itemize}
    \item minimizing sequence,
    \item weak convergence,
    \item weak lower semicontinuity.
	\end{itemize}

\HLINE


We want to show that the functional

\begin{align*}
J(u)=\int_U |\nabla u|^2\,dx-\int_U fu\,dx
\end{align*}
attains its minimum on $H_0^1(U)$, i.e. that there exists $u_0\in H_0^1(U)$ such that
\begin{align*}
J(u_0)=\inf_{u\in H_0^1(U)}J(u).
\end{align*}
This is the standard \DEFINE{direct method in the calculus of variations}, one of the foundational techniques in Calculus of Variations.
\begin{description}

\item \textbf{Step 1: Define the infimum}

Let
\begin{align*}
m=\inf_{u\in H_0^1(U)}J(u).
\end{align*}
Choose a minimizing sequence $(u_j)\subset H_0^1(U)$ such that
\begin{align*}
J(u_j)\to m.
\end{align*}
So
\begin{align*}
\int_U |\nabla u_j|^2\,dx-\int_U fu_j\,dx\to m.
\end{align*}
We want to show $(u_j)$ has a weakly convergent subsequence.

\item \textbf{Step 2: Show the minimizing sequence is bounded}

We estimate the linear term using Cauchy–Schwarz:
\begin{align*}
\left|\int_U fu_j\,dx\right| \le \|f\|_{L^2(U)}\|u_j\|_{L^2(U)}.
\end{align*}
Using Poincaré:
\begin{align*}
\|u_j\|_{L^2(U)} \le C\|\nabla u_j\|_{L^2(U)}.
\end{align*}
Hence
\begin{align*}
\left|\int_U fu_j\,dx\right| \le C\|f\|_{L^2(U)}\,\|\nabla u_j\|_{L^2(U)}.
\end{align*}
Thus
\begin{align*}
J(u_j) \ge
\|\nabla u_j\|_{L^2}^2 - C\|f\|_{L^2}\,\|\nabla u_j\|_{L^2}.
\end{align*}
Let
\begin{align*}
t_j=\|\nabla u_j\|_{L^2}.
\end{align*}
Then
\begin{align*}
J(u_j)\ge t_j^2-Ct_j.
\end{align*}
Since $J(u_j)\to m$, the values $J(u_j)$ are bounded above, so $t_j^2-Ct_j$ is bounded above, which implies $t_j$ is bounded.

Therefore,
\begin{align*}
(u_j)\text{ is bounded in }H_0^1(U).
\end{align*}

\item \textbf{Step 3: Extract a weakly convergent subsequence}

Since $H_0^1(U)$ is reflexive, there exists a subsequence (still denoted $u_j$) and some $u_0\in H_0^1(U)$ such that
\begin{align*}
u_j\rightharpoonup u_0 \qquad\text{weakly in }H_0^1(U).
\end{align*}

\item \textbf{Step 4: Use weak lower semicontinuity}

The map
\begin{align*}
u\mapsto \int_U |\nabla u|^2\,dx
\end{align*}
is weakly lower semicontinuous, so
\begin{align*}
\int_U |\nabla u_0|^2\,dx \le \liminf_{j\to\infty} \int_U |\nabla u_j|^2\,dx.
\end{align*}
Also, the linear functional
\begin{align*}
u\mapsto \int_U fu\,dx
\end{align*}
is weakly continuous, so
\begin{align*}
\int_U fu_j\,dx \to \int_U fu_0\,dx.
\end{align*}
Therefore,
\begin{align*}
J(u_0) = \int_U |\nabla u_0|^2\,dx-\int_U fu_0\,dx
\end{align*}
satisfies
\begin{align*}
J(u_0) \le \liminf_{j\to\infty} J(u_j).
\end{align*}
Since $J(u_j)\to m$,
\begin{align*}
J(u_0)\le m.
\end{align*}
But $m$ is the infimum, so
\begin{align*}
J(u_0)=m.
\end{align*}
Thus the infimum is attained.

\item \textbf{Final conclusion}

There exists $u_0\in H_0^1(U)$ such that
\begin{align*}
\boxed{
J(u_0)= \inf_{u\in H_0^1(U)} \left( \int_U |\nabla u|^2\,dx-\int_U fu\,dx \right).
}
\end{align*}

So the minimization problem has a minimizer.

\item \textbf{Why this works}

The direct method always uses:

		\begin{itemize}
        \item 1. Minimize:

Take a minimizing sequence.

        \item 2. Compactness:

Show boundedness, then extract weak convergence.

        \item 3. Lower semicontinuity:

Pass to the limit.

		\end{itemize}

This gives existence of a minimizer.

\item \textbf{PDE interpretation}

The minimizer $u_0$ is the weak solution of
\begin{align*}
-\Delta u=\frac{f}{2}
\end{align*}
(up to constants in the functional).

Thus the variational problem corresponds to an elliptic PDE.

This variational viewpoint is central in Partial Differential Equations.

---
\end{description}

If you'd like, next I can show how to derive the \DEFINE{Euler–Lagrange equation} of this minimization problem and recover the PDE satisfied by the minimizer.


\newpage\item \textbf{9. Fundamental solution of the heat equation}

Verify that
\begin{align*}
\Phi(x,t)=\frac{1}{(4\pi t)^{n/2}}e^{-|x|^2/4t}
\end{align*}
solves the heat equation.

Skills tested:

	\begin{itemize}
    \item direct differentiation,
    \item fundamental solutions.
	\end{itemize}

\newpage\item \textbf{10. Fourier transform solution of the heat equation}

Solve
\begin{align*}
u_t-\Delta u=0,\qquad u(x,0)=g(x)
\end{align*}
on $\mathbb R^n$ using the Fourier transform.

Show:
\begin{align*}
u(x,t)=\Phi(\cdot,t)\ast g(x).
\end{align*}

Skills tested:
	\begin{itemize}
    \item Fourier transform,
    \item ODE in frequency space,
    \item convolution.
	\end{itemize}

\newpage\item \textbf{ Why these 10 matter}

These ten problems cover the most common exam strategies:

	\begin{description}
    \item \DEFINE{Energy methods}

Problems 1,2,7

These test uniqueness and estimates.

    \item \DEFINE{Weak solutions}

Problems 3,4,6

These test Sobolev space fluency.

    \item \DEFINE{Sobolev/compactness}

Problems 5,6,8

These test existence methods.

    \item \DEFINE{Explicit formulas}

Problems 9,10

These test computational PDE techniques.
	\end{description}

\item \textbf{If I were prioritizing for exams:}

Study these in order:
\begin{enumerate}
\item \DEFINE{Energy estimates}
\item \DEFINE{Weak formulations}
\item \DEFINE{Lax–Milgram}
\item \DEFINE{Sobolev embeddings}
\item \DEFINE{Compactness arguments}
\item \DEFINE{Maximum principle}
\item \DEFINE{Fourier methods}
\end{enumerate}

That sequence covers a huge percentage of PDE comprehensive exams.
\end{description}

If you'd like, next I can create a \DEFINE{mock PDE comprehensive exam (3 problems, 3 hours)} at the same level as a graduate qualifying exam.


\end{document}
